\documentclass[11pt]{article}
\usepackage{mathblatt}
% Präambel des Setzers (werkzeuge/setzer-praeambel.tex, wortgleich); dazu \gteil (Muster T2U)
\makeatletter
\setlength{\mbpfbe}{0mm}% keine Punkte (pruefung.md § 5: nur beim Prüfstein)
\newlength{\gEin}\setlength{\gEin}{0pt}
\let\gSchrift\relax
\renewcommand{\pfkreuzzeile}[1]{\par\noindent\hangindent3.2em\hangafter1\hspace*{1.6em}$\square$\ #1\par}
\newcommand{\gkreuz}[1]{$\square$\ #1\hspace{2em}}
\newcommand{\gkreuzzeile}[2]{\par\noindent\hangindent2.2em\hangafter1\mbox{#1\,$\square$}\ #2\par}% Buchstabe, Option: eine je Zeile (neben Skizze)
% links, was man ansieht, rechts, was man tut; Spaltengrenze überall an derselben Stelle (gemeinsam.md „Satz“)
\newsavebox{\gbox}\newlength{\gLinks}
% v0.10 (Lehrer 13.10.): keine starre Halbierung – die linke Spalte ist so breit wie die Figur (+ 6 mm), die
% Antwort steht direkt daneben; ist die Figur breiter als 0,62 Zeile, steht die Antwort darunter
\newcommand{\gzwei}[2]{\par\smallskip\noindent\sbox{\gbox}{#1}%
  \setlength{\gLinks}{\dimexpr\wd\gbox+6mm\relax}%
  \ifdim\gLinks>0.62\linewidth
    \usebox{\gbox}\par\smallskip\noindent\begin{minipage}{\linewidth}#2\end{minipage}%
  \else
    \begin{minipage}[t]{\gLinks}\vspace{0pt}\usebox{\gbox}\end{minipage}%
    \begin{minipage}[t]{\dimexpr\linewidth-\gLinks\relax}\vspace{0pt}\raggedright #2\end{minipage}%
  \fi\par}
% v0.9 (12.10., pruefung.md § 5 „Marke mit B, Bedingungen rechts“): drittes Argument von pfaufg = Bedingungen
% der Prüfung (v0.10: \pfohnetr = TR durchgestrichen, „FOR“), klein grau im rechten Rand auf der Grundlinie der Nummer, Zeilen mit \\
\newcommand{\gbedrand}[2]{\ifblank{#2}{}{\rlap{\hspace*{\dimexpr#1+2mm\relax}%
  \parbox[t]{15mm}{\raggedright\scriptsize\leavevmode\color{mbgrau}\strut#2}}}}
% Teil eines Bündels (a, b …): Marke links im Rand, Bedingungen rechts im Rand, auf der Zeile des Teils
\newcommand{\gteilrand}[2]{\llap{\mbpfrandmarke{#1}\hspace*{\mbpfnr}}\gbedrand{\linewidth}{#2}}
% Teil eines Bündels: Buchstabe in eigener Spalte, Inhalt eingerückt (\gEin = Einzug, damit die Spaltengrenze von
% \gzwei an derselben Stelle bleibt)
\newenvironment{gbteil}[3]{\par\smallskip\noindent\gteilrand{#2}{#3}\makebox[\mbpfteil][l]{\textbf{#1}}%
  \begin{minipage}[t]{\dimexpr\linewidth-\mbpfteil\relax}\setlength{\gEin}{\mbpfteil}\raggedright\strut\ignorespaces}%
 {\end{minipage}\par}
% Grafik höchstens so breit wie die Spalte (Bündel nebeneinander, Koordinatensystem ohne Karo)
\newsavebox{\gpbox}
\newcommand{\gpassend}[1]{\sbox{\gpbox}{#1}\ifdim\wd\gpbox>\linewidth\typeout{GPASSEND: verkleinert \the\wd\gpbox\space auf \the\linewidth}\resizebox{\linewidth}{!}{\usebox{\gpbox}}\else\usebox{\gpbox}\fi}
% Skizzenfeld (v0.10, Lehrer 13.10.): Karo ab dem linken Rand über die Zeile, „Skizze“ klein grau oben links im Feld
% (wie \pfskizzenfeld der Vorlage); der rechte Rand bleibt den Bedingungen
\newcommand{\gskizze}[1]{\par\smallskip\noindent\begin{tikzpicture}
  \clip (0,0) rectangle ({\linewidth-1mm},{#1*10mm});
  \draw[black!16,line width=0.3pt,step=5mm] (0,0) grid ({\linewidth},{#1*10mm});
  \node[anchor=north west,inner sep=2pt,font=\scriptsize,text=mbgrau,fill=white] at (0,{#1*10mm}) {Skizze};
  \end{tikzpicture}\par}
\RenewDocumentEnvironment{pfaufg}{m m m}{%
  \begin{lrbox}{\mbaufgabebox}\begin{minipage}[t]{\linewidth}%
  \gSchrift\noindent\hspace*{\gEin}\mbpfrandmarke{#1}\gbedrand{\linewidth-\gEin}{#3}%
  \begin{minipage}[t]{\mbpfnr}\strut\textbf{#2}\end{minipage}%
  \begin{minipage}[t]{\dimexpr\linewidth-\gEin-\mbpfnr\relax}\raggedright\strut\ignorespaces}%
 {\end{minipage}%
  \end{minipage}\end{lrbox}%
  \setlength{\mbaufgabehoehe}{\dimexpr\ht\mbaufgabebox+\dp\mbaufgabebox\relax}%
  \ifdefempty{\gkopfnach}{\par\addvspace{5pt}\Needspace*{\mbaufgabehoehe}}{%
    % Stufenkopf, der mit dieser (ersten) Aufgabe zusammenbleibt: Platz für beide verlangen, dann den Kopf setzen
    % (12.10., Quadratische: ein Kopf stand allein unten, die Aufgabe mit Koordinatensystem auf der nächsten Seite)
    \par\Needspace*{\dimexpr\mbaufgabehoehe+40pt\relax}\gkopfnach\gdef\gkopfnach{}\par\addvspace{5pt}}%
  \noindent\usebox{\mbaufgabebox}\par\addvspace{5pt}}
% Stufenkopf vormerken: \gkopfvor{<Kopf>} – die nächste pfaufg setzt ihn, sobald beide zusammen Platz haben
\def\gkopfnach{}
\newcommand{\gkopfvor}[1]{\gdef\gkopfnach{#1}}
% Teilaufgabe mit Skizze links in gzwei: Buchstabe in eigener Spalte, darunter Text und Skizze
\newcommand{\gteil}[2]{\makebox[\mbpfteil][l]{\textbf{#1}}\parbox[t]{\dimexpr0.5\textwidth-\gEin-\mbpfnr-5mm-\mbpfteil\relax}{\raggedright #2}}
% Stufenkopf der Vorlage (mathblatt.sty 2026-10-12: Linie, Abstand, 13 pt, \pfjahre) ohne ihren Vorhalt von
% 14 Zeilen: der Setzer fordert vor dem Kopf 0,16 Seitenhöhe (Kopf + erste Aufgabe), mehr nicht
\setlength{\mbpfstufevorhalt}{0pt}
\makeatother
\newcommand{\karo}[1]{\par\smallskip\noindent\begin{tikzpicture}
  \clip (0,0) rectangle ({\linewidth-1mm},{#1*10mm});
  \draw[black!16,line width=0.3pt,step=5mm] (0,0) grid ({\linewidth},{#1*10mm});
  \end{tikzpicture}\par}
\newcommand{\kk}[1]{$\square$\,#1\hspace{0.9em}}
\newcommand{\linie}[1]{\,{\color{black!45}\rule[-1pt]{#1}{0.4pt}}\,}
\newcommand{\pkt}{\textperiodcentered{}}
% v0.14 (14.10.): Stufenkopf mit fünf Punkten (Option = Jahre für die Bestellkarte) oder „zuletzt 19“
\NewDocumentCommand{\gpunkte}{o m}{#2}
\newcommand{\gzuletzt}[1]{zuletzt #1}
% v0.7 (12.10.): Schülerfrage unter dem Stufenkopf; Ankreuzoption als Bild (Buchstabe und Kästchen darunter)
\newcommand{\gfrage}[1]{\par\nobreak\noindent\hspace*{\mbpfnr}{\small\itshape #1}\par\nobreak\addvspace{2pt}}
\newlength{\gbildbreite}\setlength{\gbildbreite}{0.32\linewidth}% Setzer setzt sie je Reihe (vier Bilder: 0,24)
\newcommand{\gbildoption}[3]{\begin{minipage}[t][#3][b]{\gbildbreite}\centering #2\par\smallskip\mbox{#1\,$\square$}\end{minipage}}
% Lineare (12.10.): Bild mit Buchstabe ohne Kästchen (Diagramm zuordnen); Aussage mit „wahr / falsch“ zum Ankreuzen
\newcommand{\gbildname}[3]{\begin{minipage}[t][#3][b]{\gbildbreite}\centering #2\par\smallskip\mbox{\textbf{#1}}\end{minipage}}
\newcommand{\gwahrfalsch}[2]{\par\noindent\hangindent2.2em\hangafter1\makebox[2.2em][l]{(#1)}#2\hfill\mbox{wahr $\square$\quad falsch $\square$}\par}
\begin{document}

\pfheftstil
\pflaufkopf[P10 \pkt{} Prüfung]{Lineare Funktionen}{PA10-7}\pfseitevon
\pfheftkopf{Lineare Funktionen}{}
\pfstufenzeile{Gerade zeichnen \pkt{} Gleichung ablesen \pkt{} Gleichung aus zwei Punkten \pkt{} Nullstelle gesucht \pkt{} Punktprobe \pkt{} Gleichung aus Text \pkt{} Gerade und Parabel}

% ===================== Stufe: Gerade zeichnen =====================
\Needspace*{0.16\textheight}\pfstufekopf{Gerade zeichnen}{\gpunkte[26 25 24 23 22]{\textcolor{black!80}{$\bullet$}\,\textcolor{black!80}{$\bullet$}\,\textcolor{black!80}{$\bullet$}\,\textcolor{black!80}{$\bullet$}\,\textcolor{black!80}{$\bullet$}}}\gfrage{Wo steht die Gleichung der linearen Funktion mit m und n? Formelsammlung, Seite~\linie{12mm}}
% Schülerfrage der Stufe (Gliederung; Einstieg nur von dort, nie aus der Bank – 12.10.)
\begin{pfaufg}{}{1.}{}
Welche Funktionen sind linear? Kreuze an. Wähl eine lineare aus und gib ihr m und n an: $m = \underline{\hspace{10mm}}$, $n = \underline{\hspace{10mm}}$.
\par\smallskip
\par\noindent\makebox[\mbpfteil][l]{\textbf{a)}}$f(x) = 3 - 2x$\hfill\kk{linear}
\par\noindent\makebox[\mbpfteil][l]{\textbf{b)}}$f(x) = x^2 + 1$\hfill\kk{linear}
\par\noindent\makebox[\mbpfteil][l]{\textbf{c)}}$f(x) = 0{,}5x$\hfill\kk{linear}
\par\fusshilfe{\mbox{1:~a) und c)}}
\end{pfaufg}
% 2025-OS-B1f-m1 (Original 2025-OS-B1f, Schwere I, rang 1; herausgelöste Fassung: keine)
\begin{pfaufg}{’25~\pkt{}~1f}{2.}{Basis}
Kreuze das Bild an, das den Graphen einer linearen Funktion zeigt.
\par\smallskip\noindent\setlength{\gbildbreite}{\dimexpr0.96\linewidth/4\relax}\sbox{\gbox}{\vbox{\begin{tikzpicture}[x=1.1cm,y=1.1cm,line width=0.6pt]\draw[-{Stealth[length=1.5mm]}] (-0.8,0)--(0.95,0);\draw[-{Stealth[length=1.5mm]}] (0,-0.6)--(0,1.05);\draw[line width=0.8pt] (-0.8,0.94)--(0.9,-0.42);\end{tikzpicture}\par\begin{tikzpicture}[x=1.1cm,y=1.1cm,line width=0.6pt]\draw[-{Stealth[length=1.5mm]}] (-0.8,0)--(0.95,0);\draw[-{Stealth[length=1.5mm]}] (0,-0.6)--(0,1.05);\draw[line width=0.8pt,domain=0.1:0.9,samples=40,smooth] plot (\x,{0.09/\x});\end{tikzpicture}\par\begin{tikzpicture}[x=1.1cm,y=1.1cm,line width=0.6pt]\draw[-{Stealth[length=1.5mm]}] (-0.8,0)--(0.95,0);\draw[-{Stealth[length=1.5mm]}] (0,-0.6)--(0,1.05);\draw[line width=0.8pt,domain=-0.6:0.6,samples=40,smooth] plot (\x,{3*\x*\x-0.4});\end{tikzpicture}\par\begin{tikzpicture}[x=1.1cm,y=1.1cm,line width=0.6pt]\draw[-{Stealth[length=1.5mm]}] (-0.8,0)--(0.95,0);\draw[-{Stealth[length=1.5mm]}] (0,-0.6)--(0,1.05);\draw[line width=0.8pt] (0.1,-0.3)--(0.5,0.5)--(0.1,0.7);\end{tikzpicture}}}\gbildoption{A}{\begin{tikzpicture}[x=1.1cm,y=1.1cm,line width=0.6pt]\draw[-{Stealth[length=1.5mm]}] (-0.8,0)--(0.95,0);\draw[-{Stealth[length=1.5mm]}] (0,-0.6)--(0,1.05);\draw[line width=0.8pt] (-0.8,0.94)--(0.9,-0.42);\end{tikzpicture}}{\dimexpr\ht\gbox/4+\baselineskip\relax}\hfill\gbildoption{B}{\begin{tikzpicture}[x=1.1cm,y=1.1cm,line width=0.6pt]\draw[-{Stealth[length=1.5mm]}] (-0.8,0)--(0.95,0);\draw[-{Stealth[length=1.5mm]}] (0,-0.6)--(0,1.05);\draw[line width=0.8pt,domain=0.1:0.9,samples=40,smooth] plot (\x,{0.09/\x});\end{tikzpicture}}{\dimexpr\ht\gbox/4+\baselineskip\relax}\hfill\gbildoption{C}{\begin{tikzpicture}[x=1.1cm,y=1.1cm,line width=0.6pt]\draw[-{Stealth[length=1.5mm]}] (-0.8,0)--(0.95,0);\draw[-{Stealth[length=1.5mm]}] (0,-0.6)--(0,1.05);\draw[line width=0.8pt,domain=-0.6:0.6,samples=40,smooth] plot (\x,{3*\x*\x-0.4});\end{tikzpicture}}{\dimexpr\ht\gbox/4+\baselineskip\relax}\hfill\gbildoption{D}{\begin{tikzpicture}[x=1.1cm,y=1.1cm,line width=0.6pt]\draw[-{Stealth[length=1.5mm]}] (-0.8,0)--(0.95,0);\draw[-{Stealth[length=1.5mm]}] (0,-0.6)--(0,1.05);\draw[line width=0.8pt] (0.1,-0.3)--(0.5,0.5)--(0.1,0.7);\end{tikzpicture}}{\dimexpr\ht\gbox/4+\baselineskip\relax}\par
\fusshilfe{\mbox{2:~A}}
\end{pfaufg}
% 2024-OS-B1i-m1 (Original 2024-OS-B1i, Schwere I, rang 2; herausgelöste Fassung: keine)
\begin{pfaufg}{’24~\pkt{}~1i}{3.}{Basis}
Gegeben ist die Gerade $y = 1{,}5x + 1$. Zeichne den Punkt ein, an dem sie die y-Achse trifft.
\par\smallskip\noindent\begin{ksys}[xmin=-3,xmax=3,ymin=-3,ymax=3,karo=0.5]\end{ksys}\par
\fusshilfe{\mbox{3:~Punkt (0|1)}}
\end{pfaufg}
% 2024-OS-K3a-m1 (Original 2024-OS-K3a, Schwere I, rang 3; herausgelöste Fassung: keine)
\begin{pfaufg}{nach~P10\\’24~\pkt{}~3a}{4.}{}
Die lineare Funktion $f$ hat die Gleichung $f(x) = 2x + 3$. Zeichne ihren Graphen in das Koordinatensystem.
\par\smallskip\noindent\gpassend{\begin{ksys}[xmin=-4,xmax=4,ymin=-5,ymax=8,karo=0.5]\end{ksys}}\par
\fusshilfe{\mbox{4:~Gerade durch (0|3) und (1|5)}}
\end{pfaufg}

% ===================== Stufe: Gleichung aus Graph =====================
\gkopfvor{\pfstufekopf{Gleichung ablesen}{\gzuletzt{19}}}
% kein Einstieg: keine Schülerfrage in der Gliederung (Schülerfrage, -Fälle, -Form, -Weg, -Skizze; oder Frage) – Einstieg nur von dort, nie aus der Bank (pruefung.md § 5, 12.10.)
% Teile eines Originals (2019-OS-K2): 2019-OS-K2b-m1, 2019-OS-K2c-m1
\begin{pfaufg}{}{5.}{}
In einem Koordinatensystem sind die Graphen $f$, $g$ und $h$ von drei linearen Funktionen gezeichnet.
% 2019-OS-K2b-m1 (Original 2019-OS-K2b, Schwere I, rang 11; herausgelöste Fassung: keine)
\begin{gbteil}{a)}{\raisebox{\dimexpr-\height+\ht\strutbox\relax}{\shortstack[r]{nach~P10\\’19~\pkt{}~2b}}\hspace*{\mbpfnr}}{}
Schreib zu jeder Eigenschaft den Graphen ($f$, $g$ oder $h$), der sie hat: „Der Graph hat die Steigung $-3$.“ „Der Graph liegt parallel zur x-Achse.“
\gzwei{\begin{ksys}[xmin=-5,xmax=5,ymin=-5,ymax=5,karo=0.5]\gerade[0.6]{-3}{3}{f}\gerade[1]{1}{3}{g}\gerade[3.5]{0}{-1}{h}\end{ksys}}{\vspace{2mm}\par\noindent Steigung $-$3: \linie{30mm}\hphantom{\,cm²}\par\smallskip\par\noindent parallel zur x-Achse: \linie{30mm}\hphantom{\,cm²}}
\end{gbteil}
% 2019-OS-K2c-m1 (Original 2019-OS-K2c, Schwere II, rang 13; herausgelöste Fassung: keine)
\begin{gbteil}{b)}{\raisebox{\dimexpr-\height+\ht\strutbox\relax}{\shortstack[r]{nach~P10\\’19~\pkt{}~2c}}\hspace*{\mbpfnr}}{}
Ordne jedem Graphen $f$, $g$ und $h$ eine Gleichung aus dieser Liste zu: $y = -3x + 3$, $y = -\tfrac{1}{3}x + 3$, $y = -1$, $y = x + 3$, $y = -3x$, $y = x - 3$.
\gzwei{\begin{ksys}[xmin=-5,xmax=5,ymin=-5,ymax=5,karo=0.5]\gerade[0.6]{-3}{3}{f}\gerade[1]{1}{3}{g}\gerade[3.5]{0}{-1}{h}\end{ksys}}{\vspace{2mm}\par\noindent f: \linie{30mm}\hphantom{\,cm²}\par\smallskip\par\noindent g: \linie{30mm}\hphantom{\,cm²}\par\smallskip\par\noindent h: \linie{30mm}\hphantom{\,cm²}}
\end{gbteil}
\fusshilfe{\mbox{5:~a)~h; b)~$-$1}}
\end{pfaufg}

% ===================== Stufe: Gleichung aus zwei Punkten =====================
\gkopfvor{\pfstufekopf{Gleichung aus zwei Punkten}{\gpunkte[25]{\textcolor{black!25}{$\circ$}\,\textcolor{black!80}{$\bullet$}\,\textcolor{black!25}{$\circ$}\,\textcolor{black!25}{$\circ$}\,\textcolor{black!25}{$\circ$}}}}
% kein Einstieg: keine Schülerfrage in der Gliederung (Schülerfrage, -Fälle, -Form, -Weg, -Skizze; oder Frage) – Einstieg nur von dort, nie aus der Bank (pruefung.md § 5, 12.10.)
% 2025-OS-K5a-m1 (Original 2025-OS-K5a, Schwere II, rang 6; herausgelöste Fassung: keine)
\begin{pfaufg}{nach~P10\\’25~\pkt{}~5a}{6.}{}
Der Graph einer linearen Funktion $f$ geht durch die Punkte $A(-2\,|\,5)$ und $B(3\,|\,-2{,}5)$. Zeichne den Graphen von $f$ in das Koordinatensystem. Entscheide, ob die Aussagen wahr oder falsch sind:
\par\smallskip\noindent\gpassend{\begin{ksys}[xmin=-4,xmax=4,ymin=-4,ymax=8,karo=0.5]\end{ksys}}\par
\gwahrfalsch{1}{Der Graph von $f$ steigt monoton.}\gwahrfalsch{2}{Der Graph von $f$ schneidet die y-Achse in $(0\,|\,2)$}
\par\smallskip\noindent Gib eine Gleichung von $f$ an.
\par\noindent\hfill y $=$\linie{40mm}
\fusshilfe{\mbox{6:~Gerade durch A und B; Aussage 1 falsch, Aussage 2 wahr; $-$1,5x + 2}}
\end{pfaufg}

% ===================== Stufe: Nullstelle =====================
\gkopfvor{\pfstufekopf{Nullstelle gesucht}{\gpunkte[22]{\textcolor{black!25}{$\circ$}\,\textcolor{black!25}{$\circ$}\,\textcolor{black!25}{$\circ$}\,\textcolor{black!25}{$\circ$}\,\textcolor{black!80}{$\bullet$}}}}
% kein Einstieg: keine Schülerfrage in der Gliederung (Schülerfrage, -Fälle, -Form, -Weg, -Skizze; oder Frage) – Einstieg nur von dort, nie aus der Bank (pruefung.md § 5, 12.10.)
% 2022-OS-K3a-m1 (Original 2022-OS-K3a, Schwere I, rang 7; herausgelöste Fassung: keine)
\begin{pfaufg}{nach~P10\\’22~\pkt{}~3a}{7.}{}
Gegeben ist die lineare Funktion $f(x) = -2x + 5$. Zeichne den Graphen von $f$ in das Koordinatensystem. Gib die Nullstelle von $f$ an.
\par\smallskip\noindent\gpassend{\begin{ksys}[xmin=-4,xmax=5,ymin=-4,ymax=7,karo=0.5]\end{ksys}}\par
\par\noindent\hfill x $=$\linie{30mm}
\fusshilfe{\mbox{7:~Gerade durch (0|5) und (1|3); 2,5}}
\end{pfaufg}

% ===================== Stufe: Punktprobe =====================
\gkopfvor{\pfstufekopf{Punktprobe}{\gpunkte[26 25 24 23 22]{\textcolor{black!80}{$\bullet$}\,\textcolor{black!80}{$\bullet$}\,\textcolor{black!80}{$\bullet$}\,\textcolor{black!80}{$\bullet$}\,\textcolor{black!80}{$\bullet$}}}}
% kein Einstieg: keine Schülerfrage in der Gliederung (Schülerfrage, -Fälle, -Form, -Weg, -Skizze; oder Frage) – Einstieg nur von dort, nie aus der Bank (pruefung.md § 5, 12.10.)
% 2023-OS-B1i-m1 (Original 2023-OS-B1i, Schwere I, rang 8; herausgelöste Fassung: keine)
\begin{pfaufg}{’23~\pkt{}~1i}{8.}{Basis}
Die Gerade $f$ hat die Gleichung $f(x) = -5x + 2$. Einer dieser Punkte liegt nicht auf $f$. Kreuze ihn an:
\par\smallskip\noindent\pfkreuz{A}{$(-2\,|\,12)$}\pfkreuz{B}{$(3\,|\,-12)$}\pfkreuz{C}{$(4\,|\,-18)$.}
\fusshilfe{\mbox{8:~B}}
\end{pfaufg}
% 2026-FOR-K5b-m1 (Original 2026-FOR-K5b, Schwere I, rang 10; herausgelöste Fassung: keine)
\begin{pfaufg}{nach~P10\\’26~\pkt{}~5b}{9.}{}
Prüfe mit einer Rechnung, ob der Punkt $P(-2\,|\,7)$ auf dem Graphen von $f(x) = -3x + 1$ liegt.
\karo{3}
\fusshilfe{\mbox{9:~P liegt auf dem Graphen}}
\end{pfaufg}

% ===================== Stufe: Gleichung aus Text =====================
\Needspace*{0.16\textheight}\pfstufekopf{Gleichung aus Text}{\gpunkte[23 22]{\textcolor{black!25}{$\circ$}\,\textcolor{black!25}{$\circ$}\,\textcolor{black!25}{$\circ$}\,\textcolor{black!80}{$\bullet$}\,\textcolor{black!80}{$\bullet$}}}
% Schülerfrage der Stufe (Gliederung; Einstieg nur von dort, nie aus der Bank – 12.10.)
\begin{pfaufg}{}{10.}{}
Welche Zahl ist m, welche n? Rechne nicht.
\par\smallskip
\par\noindent\makebox[\mbpfteil][l]{\textbf{a)}}Handytarif: 8 € im Monat und 0,09 € je Minute\hfill\linie{30mm}
\par\noindent\makebox[\mbpfteil][l]{\textbf{b)}}Kerze: 25 cm hoch, brennt je Stunde 2 cm ab\hfill\linie{30mm}
\par\fusshilfe{\mbox{10:~a) 8 ist n, 0,09 ist m; b) 25 ist n, $-$2 ist m}}
\end{pfaufg}
% Teile eines Originals (2022-OS-K6): 2022-OS-K6a-m1, 2022-OS-K6b-m1
\begin{pfaufg}{}{11.}{}
Maxi und Paula haben zum 16. Geburtstag Geld bekommen und legen jeden Monat etwas dazu, um den Führerschein zu bezahlen. Maxi beginnt mit 1 360 € und spart 24 € im Monat. Paula beginnt mit 820 € und spart 65 € im Monat.
% 2022-OS-K6a-m1 (Original 2022-OS-K6a, Schwere I, rang 12; herausgelöste Fassung: keine)
\begin{gbteil}{a)}{\raisebox{\dimexpr-\height+\ht\strutbox\relax}{\shortstack[r]{nach~P10\\’22~\pkt{}~6a}}\hspace*{\mbpfnr}}{}
Nach einem Jahr hat Paula 1 600 € gespart. Berechne, wie viel Geld Maxi nach einem Jahr hat.
\karo{2}
\end{gbteil}
% 2022-OS-K6b-m1 (Original 2022-OS-K6b, Schwere II, rang 13; herausgelöste Fassung: keine)
\begin{gbteil}{b)}{\raisebox{\dimexpr-\height+\ht\strutbox\relax}{\shortstack[r]{nach~P10\\’22~\pkt{}~6b}}\hspace*{\mbpfnr}}{FOR}
Stelle eine Gleichung für Paulas Guthaben auf ($y$: Guthaben in €, $x$: Zahl der Monate). Für Fahrschule und Prüfung braucht Paula 1 800 €. Bestimme, wie viele Monate sie mindestens sparen muss.
\karo{4}
\end{gbteil}
\fusshilfe{\mbox{11:~a)~1648 €; b)~65x + 820; 16 Monate}}
\end{pfaufg}
% Teile eines Originals (2023-OS-K3): 2023-OS-K3a-m1, 2023-OS-K3b-m1
\begin{pfaufg}{}{12.}{}
Herr Mert braucht für seinen Umzug einen Transporter. Angebot 1: 25 € Miete je Tag, 320 km insgesamt frei, jeder weitere Kilometer kostet 0,35 €. Angebot 2: einmalig 115 € Miete für eine Woche, jeder Kilometer kostet 0,08 €.
% 2023-OS-K3a-m1 (Original 2023-OS-K3a, Schwere II, rang 14; herausgelöste Fassung: keine)
\begin{gbteil}{a)}{\raisebox{\dimexpr-\height+\ht\strutbox\relax}{\shortstack[r]{nach~P10\\’23~\pkt{}~3a}}\hspace*{\mbpfnr}}{}
Die Kosten für einen Miettag sollen für beide Angebote in einem Diagramm dargestellt werden. Gib zu jedem Angebot das passende Diagramm an.
\par\smallskip\noindent\setlength{\gbildbreite}{\dimexpr0.96\linewidth/4\relax}\sbox{\gbox}{\vbox{\begin{tikzpicture}[x=1.3cm,y=1.3cm,line width=0.6pt]\draw[-{Stealth[length=1.5mm]}] (0,0)--(1.15,0) node[below left,font=\tiny] {km};\draw[-{Stealth[length=1.5mm]}] (0,0)--(0,1.05) node[below right,font=\tiny] {€};\draw[line width=0.8pt] (0,0)--(1,0.85);\end{tikzpicture}\par\begin{tikzpicture}[x=1.3cm,y=1.3cm,line width=0.6pt]\draw[-{Stealth[length=1.5mm]}] (0,0)--(1.15,0) node[below left,font=\tiny] {km};\draw[-{Stealth[length=1.5mm]}] (0,0)--(0,1.05) node[below right,font=\tiny] {€};\draw[line width=0.8pt] (0,0.35)--(1,0.6);\end{tikzpicture}\par\begin{tikzpicture}[x=1.3cm,y=1.3cm,line width=0.6pt]\draw[-{Stealth[length=1.5mm]}] (0,0)--(1.15,0) node[below left,font=\tiny] {km};\draw[-{Stealth[length=1.5mm]}] (0,0)--(0,1.05) node[below right,font=\tiny] {€};\draw[line width=0.8pt] (0,0.3)--(0.45,0.3)--(1,0.85);\end{tikzpicture}\par\begin{tikzpicture}[x=1.3cm,y=1.3cm,line width=0.6pt]\draw[-{Stealth[length=1.5mm]}] (0,0)--(1.15,0) node[below left,font=\tiny] {km};\draw[-{Stealth[length=1.5mm]}] (0,0)--(0,1.05) node[below right,font=\tiny] {€};\draw[line width=0.8pt] (0,0.7)--(0.45,0.7)--(1,0.25);\end{tikzpicture}}}\gbildname{A}{\begin{tikzpicture}[x=1.3cm,y=1.3cm,line width=0.6pt]\draw[-{Stealth[length=1.5mm]}] (0,0)--(1.15,0) node[below left,font=\tiny] {km};\draw[-{Stealth[length=1.5mm]}] (0,0)--(0,1.05) node[below right,font=\tiny] {€};\draw[line width=0.8pt] (0,0)--(1,0.85);\end{tikzpicture}}{\dimexpr\ht\gbox/4+\baselineskip\relax}\hfill\gbildname{B}{\begin{tikzpicture}[x=1.3cm,y=1.3cm,line width=0.6pt]\draw[-{Stealth[length=1.5mm]}] (0,0)--(1.15,0) node[below left,font=\tiny] {km};\draw[-{Stealth[length=1.5mm]}] (0,0)--(0,1.05) node[below right,font=\tiny] {€};\draw[line width=0.8pt] (0,0.35)--(1,0.6);\end{tikzpicture}}{\dimexpr\ht\gbox/4+\baselineskip\relax}\hfill\gbildname{C}{\begin{tikzpicture}[x=1.3cm,y=1.3cm,line width=0.6pt]\draw[-{Stealth[length=1.5mm]}] (0,0)--(1.15,0) node[below left,font=\tiny] {km};\draw[-{Stealth[length=1.5mm]}] (0,0)--(0,1.05) node[below right,font=\tiny] {€};\draw[line width=0.8pt] (0,0.3)--(0.45,0.3)--(1,0.85);\end{tikzpicture}}{\dimexpr\ht\gbox/4+\baselineskip\relax}\hfill\gbildname{D}{\begin{tikzpicture}[x=1.3cm,y=1.3cm,line width=0.6pt]\draw[-{Stealth[length=1.5mm]}] (0,0)--(1.15,0) node[below left,font=\tiny] {km};\draw[-{Stealth[length=1.5mm]}] (0,0)--(0,1.05) node[below right,font=\tiny] {€};\draw[line width=0.8pt] (0,0.7)--(0.45,0.7)--(1,0.25);\end{tikzpicture}}{\dimexpr\ht\gbox/4+\baselineskip\relax}\par
\par\noindent\hfill Angebot 1 $=$\linie{30mm}
\par\noindent\hfill Angebot 2 $=$\linie{30mm}
\end{gbteil}
% 2023-OS-K3b-m1 (Original 2023-OS-K3b, Schwere II, rang 15; herausgelöste Fassung: keine)
\begin{gbteil}{b)}{\raisebox{\dimexpr-\height+\ht\strutbox\relax}{\shortstack[r]{nach~P10\\’23~\pkt{}~3b}}\hspace*{\mbpfnr}}{FOR}
Herr Mert mietet den Transporter von Montag bis Freitag und fährt insgesamt 450 km. Berechne die Kosten für beide Angebote. Gib an, welches billiger ist. Stelle für Angebot 2 eine Gleichung auf, mit der man die Kosten einer Woche für jede Kilometerzahl berechnen kann.
\karo{4}
\end{gbteil}
\fusshilfe{\mbox{12:~a)~C; B; b)~170,50 €; 151,00 €; Angebot 2 günstiger; 0,08x + 115}}
\end{pfaufg}

% ===================== Stufe: Gemischt: Gerade und Parabel =====================
\gkopfvor{\pfstufekopf{Gerade und Parabel}{\gpunkte[26 25 24 23 22]{\textcolor{black!80}{$\bullet$}\,\textcolor{black!80}{$\bullet$}\,\textcolor{black!80}{$\bullet$}\,\textcolor{black!80}{$\bullet$}\,\textcolor{black!80}{$\bullet$}}}}
% kein Einstieg: keine Schülerfrage in der Gliederung (Schülerfrage, -Fälle, -Form, -Weg, -Skizze; oder Frage) – Einstieg nur von dort, nie aus der Bank (pruefung.md § 5, 12.10.)
% 2025-OS-B1h-m1 (Original 2025-OS-B1h, Schwere I, rang 1; herausgelöste Fassung: keine)
\begin{pfaufg}{’25~\pkt{}~1h}{13.}{Basis}
Genau einer der Werte löst die Gleichung $x(x+4) = -3$. Kreuze ihn an:
\pfkreuzab{A}{$x = 1$}\pfkreuzab{B}{$x = 3$}\pfkreuzab{C}{$x = -2$}\pfkreuzab{D}{$x = -3$.}
\fusshilfe{\mbox{13:~D}}
\end{pfaufg}
% 2026-FOR-K5a-m1 (Original 2026-FOR-K5a, Schwere I, rang 2; herausgelöste Fassung: keine)
\begin{pfaufg}{nach~P10\\’26~\pkt{}~5a}{14.}{FOR}
Im Koordinatensystem ist die Normalparabel $p$ mit $p(x) = (x-3)^2 - 1$ gezeichnet. Die lineare Funktion $f$ hat die Gleichung $f(x) = -3x + 1$. Zeichne ihren Graphen in das Koordinatensystem.
\par\smallskip\noindent\gpassend{\begin{ksys}[xmin=-3,xmax=6,ymin=-3,ymax=6,karo=0.5]\parabel{1}{3}{-1}{p}\end{ksys}}\par
\fusshilfe{\mbox{14:~Gerade durch (0|1) und (1|$-$2)}}
\end{pfaufg}
% Teile eines Originals (2024-OS-K3): 2024-OS-K3c-m1, 2024-OS-K3d-m1
\begin{pfaufg}{}{15.}{}
Die Gerade $f$ ist der Graph von $f(x) = 2x + 3$. Die Parabel $p$ ist der Graph von $p(x) = x^2 - 5$.
% 2024-OS-K3c-m1 (Original 2024-OS-K3c, Schwere I, rang 3; herausgelöste Fassung: keine)
\begin{gbteil}{a)}{\raisebox{\dimexpr-\height+\ht\strutbox\relax}{\shortstack[r]{nach~P10\\’24~\pkt{}~3c}}\hspace*{\mbpfnr}}{}
Prüfe mit einer Rechnung, ob der Punkt $T(-3\,|\,5)$ auf der Parabel $p$ liegt.
\karo{2}
\end{gbteil}
% 2024-OS-K3d-m1 (Original 2024-OS-K3d, Schwere II, rang 4; herausgelöste Fassung: keine)
\begin{gbteil}{b)}{\raisebox{\dimexpr-\height+\ht\strutbox\relax}{\shortstack[r]{nach~P10\\’24~\pkt{}~3d}}\hspace*{\mbpfnr}}{FOR}
Die Graphen von $f$ und $p$ haben zwei Schnittpunkte. Berechne ihre Koordinaten.
\karo{2}
\par\smallskip\par\noindent \hfill S$_1$: \linie{30mm}\hphantom{\,cm²}\par\noindent \hfill S$_2$: \linie{30mm}\hphantom{\,cm²}
\end{gbteil}
\fusshilfe{\mbox{15:~a)~Nein; b)~S$_1$(4|11); S$_2$($-$2|$-$1)}}
\end{pfaufg}
% 2023-OS-K4a-m1 (Original 2023-OS-K4a, Schwere II, rang 5; herausgelöste Fassung: keine)
\begin{pfaufg}{nach~P10\\’23~\pkt{}~4a}{16.}{}
Im Koordinatensystem ist eine nach unten geöffnete Normalparabel $p$ gezeichnet. Die Gerade $f$ hat die Gleichung $f(x) = 3x - 2$. Zeichne $f$ in das Koordinatensystem. Gib die Koordinaten eines Punktes an, in dem sich die beiden Graphen schneiden: $S(\ \ \ \,|\,\ \ \ )$. Entscheide, ob die Aussagen wahr oder falsch sind:
\par\smallskip\noindent\gpassend{\begin{ksys}[xmin=-5,xmax=5,ymin=-4,ymax=6,karo=0.5]\parabel{-1}{1}{5}{p}\end{ksys}}\par
\gwahrfalsch{1}{Die Steigung von $f$ ist negativ.}\gwahrfalsch{2}{Gerade und Parabel haben zwei Schnittpunkte.}\gwahrfalsch{3}{$f$ schneidet die y-Achse in $(-2\,|\,0)$.}
\fusshilfe{\mbox{16:~Gerade durch (0|$-$2) und (1|1); S(2|4); Aussage 2 wahr, Aussagen 1 und 3 falsch}}
\end{pfaufg}
% 2022-OS-K3c-m1 (Original 2022-OS-K3c, Schwere II, rang 6; herausgelöste Fassung: keine)
\begin{pfaufg}{nach~P10\\’22~\pkt{}~3c}{17.}{FOR}
Die lineare Funktion $f$ hat die Gleichung $f(x) = -2x + 5$. Die verschobene Normalparabel $p$ hat die Gleichung $p(x) = (x-4)^2 - 6$. Bestimme die Koordinaten der Punkte, in denen die Gerade $f$ die Parabel $p$ schneidet.
\karo{2}
\par\smallskip\par\noindent \hfill S$_1$: \linie{30mm}\hphantom{\,cm²}\par\noindent \hfill S$_2$: \linie{30mm}\hphantom{\,cm²}
\fusshilfe{\mbox{17:~S$_1$(1|3); S$_2$(5|$-$5)}}
\end{pfaufg}

\pfblattende{11}{26 FOR & 25 & 24 & 23 & 22 & 21 & 20 & 19 & 18 & 17 & 16}{5a 5b & 1f 1h 5a & 1i 3a 3c 3d & 1i 3a 3b 4a & 3a 3c 6a 6b & 2a 2b 6a 6c 6d 7a &  & 2a 2b 2c &  & 5a 5b & 1c 6a 6b 6c}%
{\pfausschnitt{}{Lineare Funktionen}{}}
\end{document}